\section{Elevator pitchs y preguntas frecuentes}

En esta sección se documentan las distintas variantes de presentación rápida (Elevator Pitch) diseñadas para diferentes audiencias, así como las respuestas estructuradas para los roles clave del equipo ante consultas de inversores, tutores o clientes.

\subsection{Variantes de elevator pitch}

\subsubsection{Variante 1: Enfoque comercial e inversores}

¿Sabían que las tres principales fugas de capital en los clubes deportivos provienen de la morosidad no detectada, los ingresos no autorizados en los estadios y la falta de optimización en el alquiler de espacios compartidos?

Soy [Nombre], parte del equipo creador de SocioUnido. Para resolver este problema de raíz, desarrollamos la única herramienta del mercado que ofrece una digitalización verdaderamente integral del club.

Nuestra plataforma inteligente centraliza absolutamente toda la gestión: desde el cobro automatizado de cuotas y el control seguro de accesos mediante códigos QR, hasta la administración detallada de disciplinas deportivas y la reserva de instalaciones.

Facilitamos y modernizamos los procesos en todos los niveles, conectando a la administración, los empleados y los usuarios en un mismo ecosistema.

Nuestro lema es claro: 'Porque el club es de los socios, y la gestión es de Socio Unido'.

Hoy estamos abriendo nuestra primera ronda de capital para escalar esta solución a nivel regional, y buscamos socios estratégicos que quieran ser parte de la transformación digital del deporte.

Dicho esto queremos preguntarles ¿Quién está listo para ser parte del futuro del fútbol a nivel global?

\subsubsection{Variante 2: Enfoque técnico y académico}

Hola, soy [Nombre] y en representación de nuestro equipo les presento SocioUnido, la única plataforma integral concebida para la digitalización absoluta de instituciones deportivas.

Si bien es de público conocimiento que los clubes pierden millones anualmente debido a la morosidad y a la ineficiente gestión de sus espacios, el verdadero desafío de ingeniería detrás de esta transformación radica en la infraestructura: la alta concurrencia masiva y la nula conectividad en las puertas de los estadios.

Para resolver la gestión diaria, construimos un ecosistema digital —compuesto por un panel web y una app móvil— soportado íntegramente por una arquitectura de microservicios en la nube, lo que nos permite escalar de forma elástica e independiente ante picos de demanda.

Sin embargo, nuestra mayor innovación tecnológica se encuentra en el control de accesos. Hemos desarrollado un sistema de validación de entradas 100\% offline. Mediante el uso de criptografía y la generación dinámica de códigos, garantizamos un ingreso seguro, rápido y tolerante a caídas de red, eliminando por completo el histórico cuello de botella en los molinetes.

A continuación, los invitamos a sumergirse en la arquitectura y el funcionamiento de nuestra plataforma.

\newpage

\subsubsection{Variante 3: Enfoque a clubes / dirigentes (Pitch de ventas)}

¿Sabían que, más allá de los trámites lentos que muchas veces alejan a los socios, su institución está sufriendo una fuga de capital por tres ineficiencias silenciosas? La morosidad no detectada, los ingresos no autorizados en los eventos y la subutilización de los espacios compartidos.

Soy [Nombre], y vengo a presentarles SocioUnido. Nosotros creamos una plataforma que pone el club en el bolsillo del socio y el control financiero total en las manos de la dirigencia. 

Para resolver estos problemas de raíz, nuestra aplicación automatiza el cobro de cuotas y el seguimiento de pagos para erradicar la morosidad; centraliza la reserva de canchas y disciplinas maximizando su ocupación; y blinda sus instalaciones con nuestra mayor innovación tecnológica: un sistema de control de accesos por QR 100\% offline, garantizando un ingreso ágil, masivo y antifraude, sin depender de la conexión a internet.

SocioUnido es una solución integral. Su diseño intuitivo asegura una adopción fácil e inmediata, tanto para el personal administrativo como para los socios de todas las edades, sin necesidad de adquirir hardware costoso.

¿Están listos para frenar la fuga de ingresos, modernizar sus procesos y ofrecerle a su comunidad la experiencia de primer nivel que se merecen?

\subsubsection{Variante 4: Enfoque networking y contacto directo}

Mucho gusto [Nombre del cliente], soy [Nombre] de SocioUnido y te quiero presentar nuestra solución.

Nuestra plataforma es exactamente el sistema de digitalización integral que su club necesita hoy para dar un salto definitivo en su gestión.

Sabemos que sus principales puntos de pérdida de dinero hoy son tres: la morosidad acumulada por falta de seguimiento, los ingresos no autorizados en los estadios y la subutilización de los espacios compartidos.

Nuestro sistema lo resuelve de raíz centralizando toda la gestión.

Le damos a la administración un panel de control total potenciado con inteligencia artificial para detectar patrones de morosidad y recuperar socios deudores, al socio una aplicación móvil para que pague sus cuotas y reserve espacios en segundos. Y además, blindamos las puertas de su club con un sistema de acceso por código QR 100\% offline, que garantiza un ingreso rápido y antifraude sin depender de internet.

La digitalización y el uso de nuestra plataforma es extremadamente simple. SocioUnido es una solución integral, con un diseño muy intuitivo tanto para empleados como socios que no requiere hardware adicional.

Me gustaría dejarte mi tarjeta y que nos podamos tomar un café pronto para poder contarte más acerca de cómo podemos llevarlos al siguiente nivel de la gestión deportiva.

\newpage

\subsection{Preguntas frecuentes generales}

\subsubsection{¿Cuál es la visión a futuro del proyecto?}

Nuestra visión es convertirnos en el sistema operativo estándar para los clubes de fútbol en Argentina, transformando definitivamente la relación entre el club y el hincha mediante la innovación tecnológica.

\subsubsection{¿Cuál es su misión?}

Nuestra misión es proveer a las comisiones directivas de una herramienta inteligente y centralizada, optimizando la administración, blindando los accesos al estadio y fidelizando a los hinchas para brindar nuevas vías de monetización.

\subsubsection{¿Qué necesitan actualmente para avanzar?}

Actualmente nos encontramos en la búsqueda de instituciones que funcionen como $"$Early Adopters$"$ (usuarios pioneros) para realizar pruebas en entornos reales, validar las métricas de nuestra plataforma y demostrar empíricamente el retorno de inversión en un escenario práctico.

\subsubsection{¿Podemos hacer una prueba?}

Absolutamente. Podemos preparar una demostración (demo) completamente funcional y personalizada de la plataforma. Además, acompañamos esta demostración con herramientas interactivas, como nuestro simulador de Retorno de Inversión (ROI), para no solo mostrarles cómo funciona el software a nivel técnico, sino también para sentarnos a hablar de números concretos y evidenciar cuánto capital dejarían de perder implementando SocioUnido.

\subsubsection{¿Cómo garantizan la seguridad de los datos sensibles tanto de los socios como del club?}

La seguridad es un pilar fundamental de nuestra arquitectura. Trabajamos exclusivamente con herramientas e infraestructuras en la nube certificadas por los más altos estándares de seguridad de la industria. Todo el flujo de datos está cifrado, utilizamos protocolos de autenticación robustos (como tokens JWT) y aseguramos que tanto la información financiera del club como los datos personales de los socios estén completamente blindados contra accesos no autorizados o fugas de información.

\newpage

\subsection{Preguntas frecuentes específicas por rol}

\subsubsection{Felipe Santino Ascencio - Responsable de Control de Calidad (QA), Documentación Técnica, Desarrollo de Negocio y Preventa}

\textbf{P: ¿Cuál es tu rol en el equipo?} \\
\textbf{R:} Mi rol abarca dos frentes clave. Por un lado, custodio la calidad técnica del software asegurando que se cumplan los criterios de aceptación y mantengo el registro metodológico del proyecto. Por otro lado, lidero la estrategia comercial, la preventa y el diseño de herramientas de conversión, como los simuladores de retorno de inversión, para acercar nuestra solución a los clubes.

\subsubsection{Lautaro Gabriel Ghosn - Responsable de Arquitectura, Producto, Bases de Datos e Integraciones}

\textbf{P: ¿Cuál es tu rol en el equipo?} \\
\textbf{R:} Soy el encargado de diseñar la estructura vertebral del sistema. Lidero la definición de la arquitectura general basada en microservicios, el modelado y optimización de las bases de datos relacionales, y el diseño de las APIs. También coordino las integraciones del software con plataformas externas y hardware de terceros.

\subsubsection{Martín Guerrero - Responsable de Front-end, Aplicación Móvil, Diseño UX/UI e Identidad de Marca}

\textbf{P: ¿Cuál es tu rol en el equipo?} \\
\textbf{R:} Soy el puente entre el código y el usuario. Centralizo el desarrollo integral de las interfaces, tanto del panel administrativo web como de la aplicación móvil. Mi objetivo es garantizar una experiencia de usuario fluida, reducir al mínimo la curva de aprendizaje para el personal del club y mantener la coherencia estética e identidad visual de la marca en todo el ecosistema.

\subsubsection{Axel Zielonka - Responsable de Inteligencia Artificial, Back-end y Relaciones Institucionales}

\textbf{P: ¿Cuál es tu rol en el equipo?} \\
\textbf{R:} Coordino la lógica de negocio en el servidor y la seguridad de los servicios. Mi principal enfoque técnico es el desarrollo y calibración de los motores de inteligencia artificial: tanto el modelo predictivo de morosidad mediante aprendizaje automático (Machine Learning) como el asistente conversacional impulsado por procesamiento de lenguaje natural (PLN). A la par, gestiono los vínculos y articulaciones con los directivos de las instituciones.
